\section{Conclusions}\label{sec:conclusions}

We have presented a battery of \result{battery/scenarios} conforming- and patch-contact
problems with analytic references, classified by what each reference assumes, and run on it two
elastic-foundation laws evaluated on exact mesh intersections, two baselines from the
literature on the same geometry and six native contact models of Drake, Chrono and MuJoCo. The
battery, with its classes, its band of admissible pressure distributions, its error terms and
its rules, is the reusable product of the work; the laws are an open, engine-independent
implementation of known models. The findings are these.
\begin{enumerate}
\item Where the answer depends on the pressure distribution, no elastic-foundation method is
  outside the band of admissible distributions; \law{distributed} and Drake hydroelastic, two
  discretizations of the same pressure model, agree with the model's own distribution within
  \result{err/bearing/distributed} and \result{err/cone/hydro}, which is consistency and not
  validation. A finite-element solution shows where in the band a continuum support falls
  (Section~\ref{sec:fem}): on the Winkler cosine for the thin sleeve of \scn{bearing}, where the
  four foundation methods are within 5\,\% of it and it cannot rank them, and at the rigid-punch
  end for the thick slab of \scn{disk}, against which \law{distributed} and Drake hydroelastic
  stop 18\,\% late. That is a limit of the Winkler foundation itself on a thick support, not of
  the discretization. A point contact gives a face the friction torque of its rim, and Gonthier's
  lumping costs the 6\,\% that its radius of gyration predicts.
\item The failures of the two laws follow from their model or from their coupling, with
  predictions that the runs match: the rolling resistance of the damped foundation makes a
  rolling orbit grow by \result{err/turntable/distributed} against 77\,\% predicted; a spinning
  disk on a foundation ten times stiffer flutters at the threshold and with the growth rates of
  a linear model, Eq.~\eqref{eq:flutter}; the explicit coupling loses a light body between
  stiff contacts above the step of Eq.~\eqref{eq:stability} and holds it below; a compliant
  foundation lets a block pitch over an edge early. One normal per pair of convex pieces stops
  the convergence of \law{distributed} on a journal in a hole with clearance.
\item The closed-form Coulomb friction of a linear pressure on a triangle makes the friction
  wrench of a flat face independent of its triangulation and of the decomposition of the
  bodies. One-point and vertex quadratures lose \result{tri/disk/spin/rim-fan/64/centroid/moment}
  and \result{tri/disk/spin/rim-fan/64/vertex/moment} of the spin torque of a disk triangulated
  as a fan from its rim, whatever the number of segments, and converge at second order on
  well-shaped meshes. The gain is a mesh that is not needed, not accuracy that a quadrature
  cannot reach.
\item Each law gives the same errors on three engines (within \result{independence/worst}
  wherever a run holds), and its step series are mostly flat, so the reported errors are
  mostly those of the model and the mesh. Gonthier's friction as published, with the bristle on
  the material slip, holds the bodies that the variant of the 2018 thesis lets creep.
\item Of \result{primary/n} paired comparisons of errors, \result{primary/significant} survive
  a correction for multiplicity: \law{distributed} has about half the error of \law{gonthier}
  on classes A and B, and about a quarter of that of Chrono SMC, the latter narrowly. On the
  rigid-body references alone nothing separates the laws from the native models.
\end{enumerate}

Two things would change the weight of these conclusions. An implicit coupling, with the
Jacobians of the contact wrench, would remove the step limit that is the main practical
disadvantage of the laws. And an experimental reference, with finite-element ones for the
other scenarios of class B and for the dynamics of the runs, would turn consistency with an
assumed pressure into validation; the static one of Section~\ref{sec:fem} covers two scenarios
(Section~\ref{sec:limitations}).
